% GENERATED FILE. Edit canonical lessons, then run npm run generate:books.
% canonical-source-hash: fnv1a64:283159b9a3e7fbae
% canonical-lessons: ML-C294-kalcappatu, ML-C294-tippetti, ML-C294-talu, ML-C294-cannala, ML-C294-vati

\chapter{Point of view, Matchbox, Padlock, Chain, Stick}
\label{ch:ml-point-of-view-matchbox-padlock-chain-stick}
\hlchaptermodality{294}

\begin{chapteropening}
\textbf{By the end of this chapter:} \emph{I can say a point of view, a matchbox, a padlock, a chain and a stick.}

Complete the last lesson of chapter 294: I can say a point of view, a matchbox, a padlock, a chain and a stick.
\end{chapteropening}

\section[\texorpdfstring{\ml{കാഴ്ചപ്പാട്} (kāḻcappāṭŭ)}{കാഴ്ചപ്പാട് (kāḻcappāṭŭ)}]{\ml{കാഴ്ചപ്പാട്} (kāḻcappāṭŭ) — a point of view}

\label{lesson:ML-C294-kalcappatu}



\begin{quote}
Before the new one: say the Malayalam for a channel, then the Malayalam for a programme.
\end{quote}

\subsection*{You'll want to know: \ml{കാഴ്ചപ്പാട്}}
\textbf{\ml{കാഴ്ചപ്പാട്}} — \emph{kāḻcappāṭŭ} — \textquotedblleft{}a point of view\textquotedblright{}.

\begin{cousinweb}
\textbf{\ml{കാഴ്ചപ്പാട്}} (\emph{kāḻcappāṭŭ}) is native Dravidian: \emph{kāḻca}, \textquotedblleft{}sight\textquotedblright{}, which DEDR 1443 lists under the seeing verb, with a second part meaning \textquotedblleft{}manner\textquotedblright{}: \textquotedblleft{}the manner of seeing\textquotedblright{}.
\end{cousinweb}

\begin{grammarlens}[title={What you've built}]
One more word for this chapter.
\end{grammarlens}

\subsection*{Your turn}
\begin{itemize}
\raggedright
  \item \emph{Say it:} \emph{kāḻcappāṭŭ}
  \item \emph{Say it:} \emph{kāḻcappāṭŭ}, once more
  \item \emph{Say it:} say \emph{kāḻcappāṭŭ}
  \item \emph{Recall:} say the Malayalam for a channel, then the Malayalam for a programme, then say \emph{kāḻcappāṭŭ} again
\end{itemize}

\begin{tcolorbox}[breakable,colback=teal!4,colframe=teal!35!black,title={Before you move on}]
What does \ml{കാഴ്ചപ്പാട്} mean? (\textquotedblleft{}A point of view\textquotedblright{}.) Say it once more.
\end{tcolorbox}
\section[\texorpdfstring{\ml{തീപ്പെട്ടി} (tīppeṭṭi)}{തീപ്പെട്ടി (tīppeṭṭi)}]{\ml{തീപ്പെട്ടി} (tīppeṭṭi) — a matchbox}

\label{lesson:ML-C294-tippetti}



\begin{quote}
Before the new one: say the Malayalam for a programme, then the Malayalam for a point of view.
\end{quote}

\subsection*{You'll want to know: \ml{തീപ്പെട്ടി}}
\textbf{\ml{തീപ്പെട്ടി}} — \emph{tīppeṭṭi} — \textquotedblleft{}a matchbox\textquotedblright{}.

\begin{cousinweb}
\textbf{\ml{തീപ്പെട്ടി}} (\emph{tīppeṭṭi}) is native Dravidian, \textquotedblleft{}the fire box\textquotedblright{}: \emph{tī}, \textquotedblleft{}fire\textquotedblright{} (DEDR 3266), and \emph{peṭṭi}, \textquotedblleft{}box\textquotedblright{} (DEDR 4388). DEDR 4388 compares Sanskrit \emph{peṭaka}.
\end{cousinweb}

\begin{grammarlens}[title={What you've built}]
One more word for this chapter.
\end{grammarlens}

\subsection*{Your turn}
\begin{itemize}
\raggedright
  \item \emph{Say it:} \emph{tīppeṭṭi}
  \item \emph{Say it:} \emph{tīppeṭṭi}, once more
  \item \emph{Say it:} say \emph{tīppeṭṭi}
  \item \emph{Recall:} say the Malayalam for a programme, then the Malayalam for a point of view, then say \emph{tīppeṭṭi} again
\end{itemize}

\begin{tcolorbox}[breakable,colback=teal!4,colframe=teal!35!black,title={Before you move on}]
What does \ml{തീപ്പെട്ടി} mean? (\textquotedblleft{}A matchbox\textquotedblright{}.) Say it once more.
\end{tcolorbox}
\section[\texorpdfstring{\ml{താഴ്} (tāḻŭ)}{താഴ് (tāḻŭ)}]{\ml{താഴ്} (tāḻŭ) — a padlock}

\label{lesson:ML-C294-talu}



\begin{quote}
Before the new one: say the Malayalam for a point of view, then the Malayalam for a matchbox.
\end{quote}

\subsection*{You'll want to know: \ml{താഴ്}}
\textbf{\ml{താഴ്}} — \emph{tāḻŭ} — \textquotedblleft{}a padlock\textquotedblright{}.

\begin{cousinweb}
\textbf{\ml{താഴ്}} (\emph{tāḻŭ}) is native Dravidian, one word with Tamil \emph{tāḻ} and Kannada \emph{tāḻ} (DEDR 3179). DEDR 3179 compares Sanskrit \emph{tāḍa(ka)}.
\end{cousinweb}

\begin{grammarlens}[title={What you've built}]
One more word for this chapter.
\end{grammarlens}

\subsection*{Your turn}
\begin{itemize}
\raggedright
  \item \emph{Say it:} \emph{tāḻŭ}
  \item \emph{Say it:} \emph{tāḻŭ}, once more
  \item \emph{Say it:} say \emph{tāḻŭ}
  \item \emph{Recall:} say the Malayalam for a point of view, then the Malayalam for a matchbox, then say \emph{tāḻŭ} again
\end{itemize}

\begin{tcolorbox}[breakable,colback=teal!4,colframe=teal!35!black,title={Before you move on}]
What does \ml{താഴ്} mean? (\textquotedblleft{}A padlock\textquotedblright{}.) Say it once more.
\end{tcolorbox}
\section[\texorpdfstring{\ml{ചങ്ങല} (caṅṅala)}{ചങ്ങല (caṅṅala)}]{\ml{ചങ്ങല} (caṅṅala) — a chain}

\label{lesson:ML-C294-cannala}



\begin{quote}
Before the new one: say the Malayalam for a matchbox, then the Malayalam for a padlock.
\end{quote}

\subsection*{You'll want to know: \ml{ചങ്ങല}}
\textbf{\ml{ചങ്ങല}} — \emph{caṅṅala} — \textquotedblleft{}a chain\textquotedblright{}.

\begin{cousinweb}
\textbf{\ml{ചങ്ങല}} (\emph{caṅṅala}) is a tadbhava of Sanskrit \emph{śṛṅkhala}, \textquotedblleft{}chain\textquotedblright{} (CDIAL 12580), reshaped in the mouth.
\end{cousinweb}

\begin{grammarlens}[title={What you've built}]
One more word for this chapter.
\end{grammarlens}

\subsection*{Your turn}
\begin{itemize}
\raggedright
  \item \emph{Say it:} \emph{caṅṅala}
  \item \emph{Say it:} \emph{caṅṅala}, once more
  \item \emph{Say it:} say \emph{caṅṅala}
  \item \emph{Recall:} say the Malayalam for a matchbox, then the Malayalam for a padlock, then say \emph{caṅṅala} again
\end{itemize}

\begin{tcolorbox}[breakable,colback=teal!4,colframe=teal!35!black,title={Before you move on}]
What does \ml{ചങ്ങല} mean? (\textquotedblleft{}A chain\textquotedblright{}.) Say it once more.
\end{tcolorbox}
\section[\texorpdfstring{\ml{വടി} (vaṭi)}{വടി (vaṭi)}]{\ml{വടി} (vaṭi) — a stick}

\label{lesson:ML-C294-vati}



\begin{quote}
Before the new one: say the Malayalam for a padlock, then the Malayalam for a chain.
\end{quote}

\subsection*{You'll want to know: \ml{വടി}}
\textbf{\ml{വടി}} — \emph{vaṭi} — \textquotedblleft{}a stick\textquotedblright{}.

\begin{cousinweb}
\textbf{\ml{വടി}} (\emph{vaṭi}) is native Dravidian, one word with Tamil \emph{vaṭi}, Kannada \emph{baḍi} and Telugu \emph{baḍita} (DEDR 5224).
\end{cousinweb}

\begin{grammarlens}[title={What you've built}]
One more word for this chapter.
\end{grammarlens}

\subsection*{Your turn}
\begin{itemize}
\raggedright
  \item \emph{Say it:} \emph{vaṭi}
  \item \emph{Say it:} \emph{vaṭi}, once more
  \item \emph{Say it:} say \emph{vaṭi}
  \item \emph{Recall:} say the Malayalam for a padlock, then the Malayalam for a chain, then say \emph{vaṭi} again
\end{itemize}

\begin{tcolorbox}[breakable,colback=teal!4,colframe=teal!35!black,title={Before you move on}]
What does \ml{വടി} mean? (\textquotedblleft{}A stick\textquotedblright{}.) Say it once more.
\end{tcolorbox}
